\section{\emph{Saeculorum Rex}: The Settlement of 2008}

On 7 October 2008 the Pontifical Commission \emph{Ecclesia Dei} issued decree N. 181/2008, commonly identified from its opening words as \emph{Saeculorum Rex}. Unlike the missing 1990 and early Florence acts, the checked record is a four-page signed Latin scan hosted on a clergy domain. Its dispositive clauses define exactly what changed.

The decree first narrates the Institute's purpose and life. It describes a priestly congregation headquartered at Gricigliano; members living common life on the model of canons; solemn worship according to the extraordinary form of the Roman Rite; patrons and Salesian motto; candidates from several regions; lay oblates; and the associated group of Sister Adorers. That preamble establishes the Commission's understanding in 2008. It is not a verbatim copy of the constitutions.

The Commission records written consent from Cardinal Ennio Antonelli, then apostolic administrator of Florence, dated 25 August 2008; consultation of the prefect responsible for institutes of consecrated life and societies of apostolic life; and consideration of judgments from many bishops. Acting under faculties from Benedict XVI, it erects the Institute as a society of apostolic life of pontifical right. It approves and confirms the constitutions \emph{ad experimentum} for five years. It names Wach prior general for a renewable six-year term.

\statusframe{Pontifical Commission \emph{Ecclesia Dei}, \emph{Saeculorum Rex}, N. 181/2008, 7 October 2008, especially pp. 3--4.}{The Institute moved from its earlier diocesan setting into a pontifical-right society; experimental constitutions and renewable general leadership gave it a Roman-supervised framework for international life.}{The act's operative formula does not itself use ``clerical,'' make the prior general a diocesan bishop, or exempt houses and public apostolates from local episcopal authority.}

\subsection{Erection, approval, and reception}

The act did not create the Institute from nothing. It recognizes founders, existing members, houses, formation, oblates, and missions. Nor was it merely a certificate of praise. Pontifical-right erection changed the competent authority and gave the international body a stable juridical identity beyond a single diocese. Experimental constitutional approval imposed a five-year horizon for testing and revision rather than definitive immutability.

The consultations named in the decree show that Roman and episcopal reception were part of the act. Florence's consent mattered because the center stood there. The consecrated-life dicastery's perspective mattered because the chosen canonical form belonged to its field, even though PCED issued the decree. Reports from bishops mattered because the society already served across dioceses. Pontifical right consolidated a network built through local reception; it did not replace that network.

\subsection{The companion decree for the Sisters}

N. 181/2008 expressly refuses to collapse the women into the male society. It says the associated group remains on the path toward society-of-apostolic-life status and notes a separate decree. The companion N. 182/2008 erects the Sister Adorers as a public association of the faithful experimentally for three years and approves their constitutions for the same period.

The two different terms---five years for the Institute's constitutions and three for the women's association---are useful evidence of distinct evaluation and status. A celebratory institutional page used expansive language about both decrees, but the operative Latin instruments control. Through the cutoff no checked public act showed the women completing the path described in 2008.

\subsection{What the decree says about liturgy}

The decree was issued after \emph{Summorum Pontificum} and uses that act's ``extraordinary form'' vocabulary. It commends solemnity and tradition and makes worship structurally central. It does not enumerate the Missal, Ritual, Pontifical, and Breviary in the way the FSSP-specific instruments do. Those differences warn against importing TPI-02's special-law history into the ICKSP.

The legal environment later changed. A decree framed by the 2007 categories can remain a major founding act while later universal law alters the terminology and regulation of public celebration. Determining the current effect of every constitutional or liturgical clause would require competent canonical interpretation. Historical continuity is not the same as a timeless exemption.

\claimcheck{``In 2008 the Holy See gave the ICKSP and Sister Adorers the same pontifical-right status.''}{N. 181/2008 erected the male Institute as a society of apostolic life of pontifical right. N. 182/2008 separately erected the women as a public association of faithful for three years, while the male decree said they remained on the path toward society status.}
